\documentclass[sigplan,screen,review,anonymous]{acmart}
\setcopyright{none}
\settopmatter{printfolios=true,printacmref=true}
\renewcommand\footnotetextcopyrightpermission[1]{}
\emergencystretch=2em
\acmConference[CGO '27]{International Symposium on Code Generation and Optimization}{2027}{Anonymous Review}
\acmYear{2027}
\acmDOI{}
\acmISBN{}
\hypersetup{pdfauthor={Anonymous},pdftitle={Obligation-Guided Reverification for Composed Compiler Pipelines},pdfsubject={Anonymous CGO 2027 research draft}}
\begin{document}
\title{Obligation-Guided Reverification for Composed Compiler Pipelines}
\author{Anonymous Author(s)}

\begin{abstract}
A compiler artifact can satisfy a generic intermediate-representation verifier while violating a semantic relation introduced by the selected language, optimization, or backend composition. Analysis invalidation prevents reuse of stale cached results, but it does not force a semantic check when no later consumer requests the invalidated fact. Always running every verifier avoids that omission at potentially unnecessary cost.

We present obligation-guided reverification. Components declare semantic facts, effects, invalidations, canonical verifier owners, and the first pipeline boundary at which each invalidation must be discharged. The scheduler either executes every due obligation or fails closed; under explicit completeness assumptions, a compilation cannot silently cross a due boundary with a registered obligation still pending. We implement the mechanism in an extensible compiler and evaluate structural-only, invalidation-only, selective-obligation, and always-verify policies through the same production verifier routes.

On a frozen boundary corpus, selective and always verification detect all 32 primary and 10 challenge operator shapes, while structural-only and invalidation-only miss 20 and four primary shapes respectively; no policy rejects any of 100 valid controls. Source-to-result studies cover 30 programs in System W and 12 programs in a public-SDK tensor language. Thirteen targeted semantic faults that escape weaker policies are rejected by selective and always verification at the declared boundary. Eight isolated ablations each remove its matching detection without rejecting its control. A separate demand-driven witness shows the key distinction: an invalidated but undemanded fact is not recomputed, whereas a due obligation is discharged independently of consumer demand. We claim relative boundary safety and earlier localization for the evaluated contracts, not general verifier completeness or whole-compilation speedup.
\end{abstract}

\begin{CCSXML}
<ccs2012>
 <concept>
  <concept_id>10011007.10011006.10011008</concept_id>
  <concept_desc>Software and its engineering~Compilers</concept_desc>
  <concept_significance>500</concept_significance>
 </concept>
 <concept>
  <concept_id>10011007.10011006.10011066</concept_id>
  <concept_desc>Software and its engineering~Software verification and validation</concept_desc>
  <concept_significance>300</concept_significance>
 </concept>
</ccs2012>
\end{CCSXML}
\ccsdesc[500]{Software and its engineering~Compilers}
\ccsdesc[300]{Software and its engineering~Software verification and validation}
\keywords{compiler correctness, extensible compilers, pass invalidation, translation validation, intermediate representation, contracts}
\maketitle

\section{Introduction}
Compiler frameworks increasingly compose frontends, dialects, optimizers, plugins, and backends around shared intermediate representations (IRs)~\cite{llvm,mlir}. This modularity scales engineering, but it breaks the inference that a passing IR verifier makes the selected compiler valid. A generic verifier can establish opcode schemas, branch targets, stack discipline, and type consistency while missing a relation introduced only by the selected composition. An optimizer may invalidate a capability fact produced by another extension; metadata may remain well formed while naming the wrong producer; or a shape-correct operation may violate the selected backend layout contract.

\paragraph{Motivating failure.}
Consider an optimizer that replaces a supported operation with a structurally legal intrinsic $q$. The generic IR admits $q$, so structural verification succeeds. The selected backend, however, lacks capability $k_q$. A conventional analysis manager marks a cached capability analysis stale, but no later pass asks for that analysis before backend lowering. Demand-driven recomputation therefore performs no check. The compiler reaches the backend with a valid IR object but an invalid selected system.

\begin{figure}[t]
\centering
\fbox{\begin{minipage}{0.94\columnwidth}
\footnotesize
\textbf{Before optimization:} IR valid; fact $backendSupports(x)$ valid.\\
$\downarrow$ optimizer emits structurally legal $q$ and invalidates the fact\\
\textbf{Lazy invalidation:} no consumer demand $\Rightarrow$ no recomputation\\
\textbf{Obligation scheduler:} pending fact reaches its due boundary $\Rightarrow$ canonical verifier runs or compilation fails closed.
\end{minipage}}
\caption{Invalidation records staleness; an obligation additionally requires discharge at a boundary even when no consumer requests the fact.}
\Description{A four-step compiler pipeline: a supported operation becomes a structurally legal but unsupported intrinsic; lazy invalidation performs no check without consumer demand; obligation-guided reverification invokes the canonical verifier at the due boundary.}
\label{fig:motivation}
\end{figure}

Pass managers already track cached analysis validity. LLVM uses \texttt{PreservedAnalyses} and transitive invalidation~\cite{llvmnpm}; MLIR similarly invalidates analyses unless a pass marks them preserved~\cite{mlirpm}. These mechanisms answer whether a result may be reused. Translation validation and verified compilation provide stronger semantic guarantees for covered transformations~\cite{pnueli98,necula00,leroy06,alive2}, but extension ecosystems often contain heterogeneous production verifiers that are not expressed in one proof system.

We study the middle layer between those approaches: invalidations create named semantic obligations with canonical owners and due boundaries. The runtime selectively discharges due obligations or executes every applicable verifier for comparison. Our contribution is not the observation that verification is useful, but a typed scheduling contract that makes omission observable and fail-closed.

This work makes four contributions:
\begin{itemize}
  \item a model of owned semantic facts, invalidations, due boundaries, and executable verifier routes;
  \item a deterministic selective scheduler with a relative boundary-safety property and explicit completeness assumptions;
  \item production-path studies comparing structural, invalidation, demand-driven, selective-obligation, and always verification;
  \item frozen evidence, source-to-result experiments in two language packages, and mechanism-level ablations that separate detection from false-positive behavior.
\end{itemize}

\section{Problem: IR Validity Is Not System Validity}
Let a selected compiler system be a tuple $S=(L,P,B,C)$ of language components, ordered passes, backend components, and their declared contracts. Let $x_i$ be the artifact at boundary $i$. A structural verifier establishes $V_{str}(x_i)$, such as legal opcodes and control-flow targets. The selected system additionally relies on semantic facts $F_i$---for example, bytecode-producer and source-node identity, verified stack effects, shape consistency, or backend capability support.

A transformation $t$ has an effect contract
\[
E_t=(R_t,P_t,I_t),
\]
where $R_t$ is the set of required facts, $P_t$ the set produced or preserved, and $I_t$ the invalidated facts. For every invalidated fact $f$, the selected contract table maps $f$ to at most one canonical verifier rule $owner(f)$. If no executable route exists, or two components claim the same canonical rule, scheduling fails closed.

The key distinction is
\[
V_{str}(x_i) \centernot\implies Valid(S,x_i).
\]
The artifact can be a valid instance of a generic IR and still violate the chosen language/pass/backend relation. A verifier route converts this gap into an executable obligation $o=(f,owner(f),i)$.

\paragraph{Capability counterexample.}
Assume the generic IR admits a structurally legal instruction $q$, the selected emission contract requires capability $k$ for $q$, and the selected backend capability set $K_B$ does not contain $k$. For an artifact $x$ containing $q$, the opcode and operand schema may satisfy $V_{str}(x)$, while the required relation $k\in K_B$ is false. Hence $V_{str}(x)$ holds but $Valid(S,x)$ does not. This witness does not depend on malformed IR; the missing information is the selected-system relation itself.

The same construction applies to identities and compiler facts: a tag or pattern may be globally declared yet attributed to the wrong producer or source node, and a transformed artifact may remain structurally legal after invalidating a fact that no later stage re-establishes. Our empirical question is therefore not whether generic structure should be checked, but which selected-system obligations must additionally be discharged and where.

\section{Scheduler and Running Example}
\subsection{Typed ownership and route resolution}
Each selected extension contributes immutable contract facets for fact ownership, pipeline effects, IR identities, and backend capabilities. Selection constructs one table for the actual runtime components. Duplicate fact owners, unknown facts, foreign production, impossible preservation, and ambiguous repeated module occurrences are rejected before semantic scheduling.

For every routable fact, the verifier registry resolves a triple $(rule,owner,b_{min})$. The rule names executable verification logic; the owner prevents an unrelated extension from satisfying the obligation with a similarly named check; and $b_{min}$ records the first boundary exposing the required artifact. Registry construction rejects conflicting owners and unknown boundaries. Effect validation rejects an invalidation that has no route, so a missing route cannot disappear before reaching the scheduler.

\subsection{Obligation creation and discharge}
When a selected transformation invalidates $f$, effect validation changes its state to invalid and emits $o=(f,rule,owner,b_c,b_e)$ with $b_e=\max(b_c,b_{min})$. The compilation observer retains the fact state and pending obligations while advancing monotonically through the boundaries. Obligations are normalized by deadline, rule, and fact. P2 selects exactly those whose first eligible boundary is current. It rejects an overdue obligation, an absent route, or an owner mismatch. P3 selects all boundary routes, but any due obligation remains attached to the corresponding invocation for lineage and diagnostics.

A successful production verifier invocation is the discharge event. It marks only facts owned by the invoked route and removes only matching due obligations; a baseline semantic check does not silently discharge a passive P1 obligation. The observer reports verifier diagnostics through the normal compiler error policy, and rejection prevents backend execution. Structural and semantic scopes are separate, avoiding a comparison in which a mandatory structural call silently performs every semantic check.

\subsection{Demand recomputation}
P1D receives an explicit set of queried facts. A query for an unknown fact fails closed. A query for an invalidated known fact schedules its canonical route. Without a query, P1D leaves the invalidation pending even if the pipeline reaches a semantic boundary. This is not an intentionally broken baseline: it models the normal guarantee of lazy analysis invalidation, where recomputation is demand-triggered rather than deadline-triggered.

\subsection{Deferred-route example}
The lifecycle is not limited to immediate post-pass checks. A production-path conformance case invalidates \texttt{backend.input-verified} at optimized AIR. Its canonical verifier requires the final AIR artifact and therefore declares backend input as its earliest executable boundary. P2 carries the pending obligation across the intervening observer return and discharges it in the newly exposed pre-backend hook. This case tests the same observer and scheduler used by normal compilation; it is an implementation-conformance witness, not an independently sampled compiler defect.

\subsection{Running example}
Figure~\ref{fig:running-example} follows one fault through the policies. A frontend produces AIR satisfying structural and selected capability facts. An optimizer performs a structurally legal rewrite to an intrinsic unsupported by the selected backend and invalidates \texttt{air.intrinsics-supported}. The generic AIR verifier still accepts the instruction schema.

P0 sees valid structure and proceeds. P1 records the invalidation but executes no semantic route. P1D rejects only in the variant where a downstream component explicitly queries the invalidated capability fact; in the matched unqueried variant it proceeds. P2 creates an obligation owned by \texttt{core.air} and executes it at optimized AIR, the first eligible boundary, before backend execution. P3 executes the same semantic verifier plus every other applicable route and yields the same correctness outcome on this case.

\begin{figure}[t]
\centering
\fbox{\begin{minipage}{0.94\columnwidth}
\small
source $\rightarrow$ frontend $\rightarrow$ valid AIR
$\rightarrow$ optimizer emits unsupported intrinsic
$\rightarrow$ optimized-AIR boundary
$\rightarrow$ backend\\[2pt]
\textbf{selected facts:} schema valid; branch targets valid; intrinsics supported\\
\textbf{effect:} invalidate intrinsics-supported\\
\textbf{P0/P1:} proceed \quad
\textbf{P1D:} reject iff queried \quad
\textbf{P2/P3:} reject at optimized AIR
\end{minipage}}
\caption{Composed pipeline and the selected contract relation in the running example. The rewrite is structurally legal; the violated relation is selected-backend capability support.}
\Description{A left-to-right compiler pipeline. An optimizer invalidates the selected backend capability fact. P0 and P1 proceed, P1D rejects only when queried, and P2 and P3 reject at the optimized-AIR boundary.}
\label{fig:running-example}
\end{figure}

\begin{figure}[t]
\centering
\fbox{\begin{minipage}{0.90\columnwidth}
\small
$\mathsf{valid}$ $\xrightarrow{invalidate}$ $\mathsf{invalid}$
$\xrightarrow{create\ o}$ $\mathsf{pending}$
$\xrightarrow[\text{sound accept}]{owner\ route}$ $\mathsf{discharged}$\\[-1pt]
\hfill $\searrow$ \textsf{rejected} (false fact, missing route, owner conflict, or missed boundary)
\end{minipage}}
\caption{Fact/obligation lifecycle. Meaning does not depend on color.}
\Description{A state transition diagram from valid fact to invalid fact, pending obligation, and either discharged or rejected.}
\label{fig:obligation-lifecycle}
\end{figure}

\section{Implementation}
We implement the model in System W, an extensible .NET compiler. The contract owner is a production module shared by selection, compilation observers, and experiment runners. A fact state exposes \texttt{valid}, \texttt{invalid}, and \texttt{unknown}; an effect contract contains normalized \texttt{requires}, \texttt{produces}, \texttt{preserves}, and \texttt{invalidates} sets; and each \texttt{VerificationObligation} stores fact, rule, canonical owner, creation boundary, and first eligible boundary.

The production observer receives actual bytecode, AIR, optimized AIR, and backend-input artifacts. The backend-input hook is part of the normal execution builder and runs immediately before backend compilation. A compilation-scoped lifecycle record is keyed by the identity of the immutable compilation input. It contains the last observed boundary, current fact state, and pending obligations. Boundary order must be strictly increasing; state is discarded on completion or failure so a failed attempt cannot contaminate a retry.

Stage seeds add only facts not already marked invalid. Effect validation then follows selected module occurrence order rather than sorting by module name. It rejects missing requirements, invalid preservation claims, repeated ambiguous occurrences, and invalidations without a canonical verifier route. A route descriptor records its earliest executable boundary. The created deadline is the later of creation and route availability, so an optimized-AIR invalidation can remain pending until backend input. The scheduler uses deterministic rule/fact ordering and rejects unknown policies, unknown demand queries, route conflicts, owner mismatches, and missed deadlines.

P0, P1, P1D, P2, and P3 share the same structural and semantic verifier implementations. P1D adds only an execution-scoped demand set; it does not alter the artifact or oracle. P2 schedules due obligation routes. P3 schedules every route exposed at the boundary. A route may own several facts, so selectivity is at the canonical-route granularity: one invocation can re-establish multiple facts. We report verifier invocations rather than presenting fact count as checker cost.

\subsection{Admissibility and boundary procedure}
Contract admission is separated from policy execution. Selection first normalizes fact identifiers, effect sets, route identifiers, owner identities, and earliest executable boundaries. It then resolves every selected fact to at most one canonical verifier owner and rejects duplicate owners, an unknown boundary, or an invalidation with no executable route. This preflight prevents a later policy choice from changing which checker is authoritative. It also keeps P1D, P2, and P3 comparable: they differ in scheduling, not in verifier implementation or route resolution.

At each observed boundary the implementation performs four ordered steps. First, it merges the stage seed into the retained state without overwriting a prior invalidation, then validates requirements and preservation claims. Second, it applies produced and invalidated effects in selected occurrence order; each invalidation creates or retains an obligation whose lineage records the triggering boundary. Third, it computes the execution set: P1D chooses queried invalid facts whose routes are executable, P2 chooses obligations whose first eligible boundary is current, and P3 chooses every canonical route exposed there. Fourth, it executes routes in deterministic rule/fact order. A successful invocation marks its owned scheduled facts valid and removes matching due obligations. Rejection, an unknown result, an owner mismatch, an unresolved route, an out-of-order boundary, an overdue obligation, or pending P2/P3 work at the final modeled boundary stops compilation.

Structural and semantic establishment are intentionally separate. AIR receives one semantic baseline when it first becomes observable. If P1 later records an invalidation, that baseline is not treated as a discharge event. Conversely, a P3 invocation with no obligation can establish the route's owned facts at its current boundary. This prevents mandatory baseline checks from making P1 appear to enforce obligation discharge.

The implementation deliberately does not infer a route from a nearby component name or from registration order. Such inference would make correctness depend on composition accidents invisible to the effect contract. Repeated occurrences are accepted only when their normalized effects and ownership are unambiguous. This is why occurrence order, owner identity, missing-route handling, conflict rejection, and deferred-boundary handling appear as production invariants and targeted tests.

\subsection{Implementation-conformance checks}
The conformance suite checks the actual orchestration path, not only the standalone scheduler. It asserts the observer order
\[
\begin{aligned}
\text{bytecode}&<\text{AIR}<\text{optimized AIR}\\
&<\text{backend input}<\text{backend compile}.
\end{aligned}
\]
A deferred case invalidates \texttt{backend.input-verified} at optimized AIR. P2 creates an obligation with creation boundary optimized AIR and first eligible boundary backend input, carries it in the compilation-scoped state, and invokes the backend-input owner before compilation. Separate tests verify carry-forward when no new effect reproduces the obligation, overdue rejection, route/owner conflicts, and state cleanup after failure. These tests establish code/model correspondence for lifecycle mechanics; they do not validate the semantic truth of arbitrary extension contracts.

\paragraph{System W language path.}
The source-to-result runner executes the normal dialect engine, optimizer, production observer, and CIL/interpreter backends in fresh processes. Research policy switches remain internal/friend APIs and do not expand the shipped facade. P3 remains the strongest checking profile, while the frozen historical experiment runner continues to count only boundaries represented by its versioned protocol.

\paragraph{Language T public-SDK path.}
A separate small language package uses only the public extension SDK. Its shape/layout facts and verifier route differ from System W's arithmetic/capability relations, while the scheduler and contract representations are shared. We do not call the package independently authored: its cases were created during this study.

\paragraph{Baseline defect repair.}
Historical source-to-result case P07 mixed double parameters with an integer literal and failed under every policy before a result. We retained the historical receipt, repaired the independent numeric-dispatch defect with an explicit finite promotion matrix, added CIL/interpreter regression coverage, versioned the corpus, and reran the unchanged source and oracle. P07 now produces 13 under every policy and is counted as a valid control only in the new protocol.

The exact pre-hardening master contract was 1,597 passing tests. New test counts are reported only after mechanical discovery on the final branch head; the historical count is not rewritten.

\section{Methodology}
\subsection{Research questions}
\begin{description}
 \item[RQ1:] Which selected-system faults escape generic structural verification?
 \item[RQ2:] What does due-boundary obligation discharge add beyond invalidation records and lazy demand-driven recomputation?
 \item[RQ3:] How do selective and always verification compare in detection, localization, false positives, and verifier work on the evaluated corpora?
 \item[RQ4:] Does the policy distinction survive source-to-result execution and a second public-SDK language package?
\end{description}

\subsection{Frozen boundary corpus}
The historical System W corpus is immutable: 40 instances representing 32 primary operator shapes in five families, 10 post-freeze challenge operators, and 100 valid controls per policy. Faults are repeated three times. The v3 study preserves historical B0/B1/B2 labels, emits a separate four-policy schema, and mechanically checks the frozen mutation catalog.

\subsection{Demand-driven witness}
To compare against a stronger invalidation baseline without changing the frozen matrix, we execute one controlled invalidation through three schedules: lazy demand with no consumer request, lazy demand with an explicit request, and P2 obligation discharge. The same AIR verifier and invalid artifact are used in the demanded and P2 cells, plus a matched valid control. The witness tests the semantic distinction, not population prevalence.

\subsection{System W source-to-result corpus}
The end-to-end corpus contains 30 programs across arithmetic, pricing, and backend strata. Five cases enable an optimizer that preserves structural validity while violating a declared semantic relation at optimized AIR. Each case/policy pair runs twice in a fresh process, yielding 240 raw records. CIL faults produce a wrong result under P0/P1; interpreter faults may reach a later backend/runtime failure. P2/P3 must reject at optimized AIR with the expected capability diagnostic. After repairing mixed numeric promotion, all 25 non-fault programs remain unchanged and must produce their expected result under every policy.

\subsection{TensorRules and ablations}
TensorRules freezes two valid programs, two intrinsically invalid programs, and eight post-verification mutations of dimensions, output shape, layout, dynamic extent, backend capability, broadcast relation, stale shape, and fact ownership. Four policies produce 48 observations. Separately, eight minimal counterexamples and matched controls isolate producer identity, source-node identity, selected pipeline order, canonical ownership, conflict rejection, unresolved obligations, backend capabilities, and repeated-occurrence rejection.

\subsection{Evidence and claim control}
Raw JSONL is written before assertions. Artifacts carry source identity, environment records, and SHA-256 manifests. Hosted CI is a reproducibility provider, not a performance environment. An external-author freeze/import kit exists, but no externally authored corpus is included in the reported results; corresponding generalization claims remain blocked.

\section{Results}
\subsection{RQ1: structural versus selected-system validity}
Table~\ref{tab:boundary} preserves the historical primary/challenge denominators and reports the new demand set separately. P0 detects 12/32 primary shapes and 1/10 challenge shapes. Passive invalidation P1 detects 28/32 and 10/10 but does not discharge the four remaining primary obligations. P2/P3 detect every evaluated primary and challenge case. No policy rejects a valid control in protocol v4.

\begin{table}[t]
\caption{Boundary protocol v4. Primary/challenge are historical denominators; demand and controls are v4 study sets.}
\label{tab:boundary}
\small
\begin{tabular}{@{}lrrrr@{}}
\toprule
Policy & Primary & Chall. & Demand & Control FP \\
\midrule
P0 structural & 12/32 & 1/10 & 0/2 & 0/120 \\
P1 invalidation & 28/32 & 10/10 & 0/2 & 0/120 \\
P1D demand & 28/32 & 10/10 & 1/2 & 0/120 \\
P2 selective & 32/32 & 10/10 & 2/2 & 0/120 \\
P3 always & 32/32 & 10/10 & 2/2 & 0/120 \\
\bottomrule
\end{tabular}
\end{table}

The primary paired difference between P2 and P0 is 20 cases; between P2 and P1/P1D it is four. These are absolute discordance counts on a fixed constructed corpus, not population estimates. P2 and P3 have no discordant evaluated boundary case.  The false-positive column is a necessary but weak check: the controls cover selected valid relations, not every legal compiler state.

\subsection{RQ2: demand versus obligation}
The matched demand cases isolate the scheduling distinction. Both invalidate the same selected capability fact and use the same production verifier route. In the queried case, P1D recomputes and rejects. In the unqueried case, P1D proceeds to a wrong result. P2 rejects both at optimized AIR because the obligation deadline is independent of consumer demand. This is the minimal executable counterexample to treating lazy invalidation as equivalent to boundary enforcement.

The result is intentionally asymmetric.  It does not show that demand-driven invalidation is unsound in general: a system whose every semantically relevant use performs a query will reject both cases.  It shows that the demand guarantee alone does not imply a boundary guarantee.  P2 trades additional contract obligations for that stronger conditional property.

\subsection{RQ3: source-to-result behavior}
The source-to-result run produced and validated 320 raw fresh-process records. All 25 valid controls, including repaired P07, return the expected result under all five policies. P1 and P0 miss all seven targeted faults. P1D rejects the one explicitly queried demand fault but otherwise exhibits the same wrong-result or late-failure behavior as P1. P2/P3 reject all seven targeted faults before backend execution. Across all 32 programs, P2 and P3 agree on classification, diagnostic family, and first detection boundary (32/32).

P07 is no longer an excluded baseline failure: the unchanged expression \texttt{x * 2 + y} with $x=5$, $y=3$ returns 13 under CIL and interpreter. Historical evidence recording the previous all-policy assertion remains immutable and is linked from the new protocol receipt.  The repair is evaluated separately from the scheduler so that accepting P07 cannot inflate any policy's fault-detection count.

P2 invokes fewer semantic routes than P3 on clean boundaries because P3 also runs routes without due obligations. Boundary-kernel tick counts are retained as qualitative work evidence only; no whole-compilation speedup is claimed.

\subsection{RQ4: public-SDK language and validity boundary}
Language T schema v2 produced 70 observations over 14 cases. P2/P3 agree on all 14 classifications. Both reject the ten post-verification faults; P1D rejects the queried demand fault and misses the matched unqueried one. The result shows that the scheduling distinction survives a different shape/layout fact vocabulary through the public SDK, but the package and cases remain study-authored and share the framework scheduler.

No human-authored frozen corpus is available. External independence therefore remains blocked rather than inferred from Language T or model-generated cases.  Historical screening found three issue-defined pre-study regressions with stable original 3/3 evidence, but the new exact-revision campaign did not start before the bounded host terminated compilation.  Consequently the historical corpus contributes provenance and feasibility evidence only, not a policy-comparison result.  The main quantitative claims remain confined to the constructed boundary, source-to-result, and Language T corpora.

\subsection{Evidence consistency}
Every reported aggregate is regenerated from versioned machine-readable records.  CI checks immutable historical catalogue hashes, validates the separate v4 demand catalogue, rejects mixed schema versions, byte-compares generated ablation tables, and runs the source-to-result and Language T hard gates.  These checks protect result lineage; they do not make the study independent or prove that the selected oracles are complete.

\section{Ablation and Verification Work}
We evaluate two complementary ablation levels. The policy-level study removes typed contracts (P0), keeps invalidation records without discharging obligations (P1), or discharges only an explicitly demanded obligation (P1D). P2 discharges every due obligation at its first eligible boundary. The mechanism-level study changes one validator behavior at a time while holding fixed one minimal counterexample and one valid control.

\begin{table}[t]
\caption{Policy-level ablation. ``Hist. W/T'' are the historical System W and TensorRules faults; ``Demand W/T'' are matched queried/unqueried demand cases. The table is regenerated from validated evidence and byte-compared in CI.}
\label{tab:policy-ablation}
\scriptsize
\centering
\begin{tabular}{@{}lrrrrr@{}}
\toprule
Variant & Prim. & Chall. & Hist. W & Hist. T & Demand W/T \\
\midrule
P0 no contracts & 12/32 & 1/10 & 0/5 & 0/8 & 0/2, 0/2 \\
P1 no discharge & 28/32 & 10/10 & 0/5 & 0/8 & 0/2, 0/2 \\
P1D demand only & 28/32 & 10/10 & 0/5 & 0/8 & 1/2, 1/2 \\
P2 selective & 32/32 & 10/10 & 5/5 & 8/8 & 2/2, 2/2 \\
\bottomrule
\end{tabular}

\end{table}

Removing typed contracts loses 20 primary and nine challenge detections relative to P2. Keeping invalidation without discharge recovers the challenge set but still loses four primary shapes and every targeted early rejection in both historical source-level studies. P1D retains the same historical behavior as P1, but rejects exactly the queried member of each matched demand pair; the unqueried member remains undischarged. P2 rejects both members because both obligations become due at their first eligible semantic boundary.

\begin{table}[t]
\caption{Eight isolated mechanism ablations. ``Full'' and ``Ablated'' report detection of the corresponding minimal counterexample; ``Control FP'' reports rejection of its valid control by the ablated validator. The table is evidence-generated and byte-compared in CI.}
\label{tab:mechanism-ablation}
\scriptsize
\centering
\begin{tabular}{@{}lp{0.43\columnwidth}ccc@{}}
\toprule
ID & Removed mechanism & Full & Ablated & Control FP \\
\midrule
M01 & Producer identity & 1/1 & 0/1 & 0/1 \\
M02 & Source identity & 1/1 & 0/1 & 0/1 \\
M03 & Selected order & 1/1 & 0/1 & 0/1 \\
M04 & Canonical owner & 1/1 & 0/1 & 0/1 \\
M05 & Route conflict & 1/1 & 0/1 & 0/1 \\
M06 & Fail-closed route & 1/1 & 0/1 & 0/1 \\
M07 & Capability contract & 1/1 & 0/1 & 0/1 \\
M08 & Repeated occurrence & 1/1 & 0/1 & 0/1 \\
\bottomrule
\end{tabular}

\end{table}

The full protocol detects all eight counterexamples. Removing each mechanism independently loses the matching detection, while neither the full nor ablated validators reject the eight controls. These are necessity witnesses for the implemented mechanisms, not estimates of natural fault frequency and not a proof that the mechanisms are jointly minimal for every compiler.

On 120 boundary controls, P2 executes 120 verifier calls and P3 executes 160, a 25\% reduction. This meets the frozen isolated-work threshold, but it is neither a pinned-machine result nor whole-compilation timing. We therefore report it only as verifier-work evidence and make no efficiency headline.

\section{Related Work}
\paragraph{Pass managers, invalidation, and verifier insertion.}
LLVM's new pass manager caches analyses and uses preservation sets to decide which cached results remain valid after a transformation~\cite{llvmnpm}. Analyses may implement transitive invalidation through their dependencies; passes may also update selected analyses manually. MLIR similarly computes analyses lazily, assumes them invalidated after a pass unless explicitly preserved, and lets an analysis inspect preservation of its dependencies~\cite{mlirpm}. These mechanisms answer whether a cached analysis result may be reused. P1D is the direct demand-driven analogue in our study: an invalidated fact is recomputed when queried.

MLIR also exposes a materially stronger operational alternative: its pass manager can run the IR verifier after each individual pass, and pass instrumentation can insert additional callbacks~\cite{mlirpm}. When one verifier expresses every required relation, this configuration can already provide an effective post-pass gate without obligation metadata. It differs from P2 in three ways: the verifier is ordinarily selected by IR kind rather than by a fact/owner route, it runs after every pass rather than only for due selected relations, and a generic IR verifier cannot enforce a relation absent from that IR's verification interface. Our contribution is therefore not the idea of checking after transformations. It is the fail-closed composition contract that identifies which selected relation became suspect, assigns exactly one production owner, records the earliest artifact boundary at which that route can execute, and carries the obligation there even without a query. We do not report a whole-compilation speed comparison with MLIR-style per-pass verification.

\paragraph{Translation validation.}
Translation validation checks each produced translation~\cite{pnueli98,necula00}; verified validators reduce trust for particular transformations~\cite{tristan08}, and Alive/Alive2 automate bounded LLVM rewrite validation~\cite{alive,alive2}. These define strong relations for specific IRs or transformations. Our complementary problem is scheduling heterogeneous existing verifiers: which selected relation became suspect, who owns it, and when it becomes mandatory. The scheduler neither replaces nor strengthens a routed validator.

\paragraph{Verified, typed, and certifying compilation.}
CompCert proves semantic preservation~\cite{leroy06,leroy09backend}; typed assembly, proof-carrying code, and certifying compilation preserve or check explicit evidence~\cite{tal,pcc,certcompiler}. These provide stronger assurance. Our conditional mechanism merely makes existing production checks unavoidable at selected composition boundaries.

\paragraph{Compiler testing and reduction.}
Csmith, EMI, and YARPGen find compiler failures~\cite{csmith,emi,yarpgen}; C-Reduce, Perses, and delta debugging reduce them~\cite{creduce,perses,delta}. Our targeted faults test scheduling, not natural prevalence or discovery power. Fuzzing can supply future external cases but does not enforce a selected relation before backend consumption.

\begin{table*}[t]
\caption{Feature comparison with the strongest material alternatives. ``Query'' means consumer-triggered execution; ``boundary'' means a declared compiler deadline. This table compares mechanisms, not assurance strength: verified compilation provides a substantially stronger whole-compiler result than our conditional theorem.}
\label{tab:related-comparison}
\scriptsize
\centering
\setlength{\tabcolsep}{2.5pt}
\begin{tabular}{@{}p{0.225\textwidth}cccccccc@{}}
\toprule
Mechanism & Inval. & Sem. facts & Canon. owner & Deadline & Missing route & Selective & S2R gate & Indep. ext. \\
\midrule
LLVM/MLIR analysis managers & yes & analysis & no & query & n/a & query & no & yes \\
Verifier after every pass & no & IR schema & IR owner & each pass & verifier failure & no & yes & yes \\
Generic verifier at chosen sites & no & schema & fixed & call site & n/a & no & no & limited \\
Translation validation / Alive2 & per transform & relation & validator & transform & n/a & per transform & usually no & relation-specific \\
Verified / certifying compiler & proof state & yes & proof/checker & proof boundary & proof failure & proof-guided & yes & proof interface \\
P1D demand recomputation & yes & selected & route owner & query & fail closed on query & query & no mandatory gate & yes \\
P2 obligation-guided & yes & selected & unique & boundary & fail closed & due only & yes & yes \\
P3 always verify & yes & selected & unique & every boundary & fail closed & all routes & yes & yes \\
\bottomrule
\end{tabular}
\end{table*}

Pass managers already invalidate lazily; per-pass verification already gates a complete IR-level checker; validators and verified compilers provide stronger relations or proofs. Our contribution is narrower: map a selected-system invalidation to one owner and deadline, carry it to the earliest executable artifact boundary, and fail closed if no route resolves.

\paragraph{Positioning.}
The strongest practical alternative combines lazy invalidation, production validators, per-pass verification, and testing. Per-pass verification closes the unqueried gap when one available checker expresses the relation. P2 targets heterogeneous or deferred routes whose required artifact or owner appears later. For a monolithic compiler with one complete always-run verifier, P2 adds little.

\section{Threats to Validity and Limitations}
\paragraph{Constructed corpora and author/model involvement.}
The primary, challenge, demand, source-to-result, and Language T faults were designed during the study and may align with implementation structure. Repetitions test reproducibility, not independent semantic sampling. We report paired case counts and all misses; we do not infer prevalence or generality. A human-authored frozen corpus remains \texttt{BLOCKED\_EXTERNAL}.

\paragraph{Historical selection.}
A historical corpus can reduce direct construction bias, but author-selected inclusion and backporting introduce their own bias. We therefore freeze criteria before search, include every eligible case, derive the oracle from the original report/fix rather than P2 output, and publish exclusions and failed replays. We do not call this corpus independent.

\paragraph{Fact and effect completeness.}
The theorem is conditional. An omitted fact, missing invalidation, or false preservation declaration can prevent obligation creation. Fail-closed missing-route handling protects only facts that are declared as invalidated; it cannot discover an unmodeled relation.

\paragraph{Verifier soundness.}
A uniquely owned route can still contain a bug. P2 and P3 share verifier implementations, so their parity cannot validate verifier soundness. The study tests selected negative and control cases but does not prove the verifier.

\paragraph{Generality.}
System W and Language T use different fact vocabularies and language paths, but share one scheduler/framework. Evaluated backends and optimizers are limited. P2/P3 parity is explicitly corpus-bounded.

\paragraph{Performance.}
Verifier-work counts and boundary-kernel timing do not establish whole-compilation performance. Shared hosted runners are not pinned. Until the machine protocol runs, overhead and speedup remain unknown and absent from the headline claims.

\paragraph{Anonymity.}
The public project and earlier research title can be found through search. Rewriting public history would be misleading. The submission instead uses a new title and a separately sanitized package, but search-based de-anonymization risk cannot be eliminated completely.

\section{Conclusion}
A generic IR may be structurally valid while violating the semantic contracts of a selected compiler composition. Obligation-guided reverification turns declared invalidations into verifier-owned duties that must be discharged by a due boundary, independent of later consumer demand. Under explicit completeness and soundness assumptions, a successful boundary transition cannot leave a registered due obligation pending.

Across a frozen System W boundary corpus, a 30-program source-to-result study, and a public-SDK TensorRules package, selective obligation discharge matches always verification on evaluated classifications and rejects faults missed by structural and invalidation-only policies. A separate lazy-demand witness isolates the key trigger difference, while eight mechanism ablations provide concrete necessity witnesses. The evidence supports relative boundary safety, detection, and localization for the evaluated contracts. It does not establish general verifier completeness, natural fault prevalence, or whole-compilation speedup; externally authored faults and pinned-machine performance remain necessary before making those claims.

\bibliographystyle{ACM-Reference-Format}\begin{thebibliography}{20}

\bibitem{llvm}
Chris Lattner and Vikram Adve.
\newblock LLVM: A compilation framework for lifelong program analysis and transformation.
\newblock In \emph{CGO}, pages 75--86, 2004. DOI: 10.1109/CGO.2004.1281665.

\bibitem{mlir}
Chris Lattner, Mehdi Amini, Uday Bondhugula, Albert Cohen, Andy Davis, Jacques Pienaar, River Riddle, Tatiana Shpeisman, Nicolas Vasilache, and Oleksandr Zinenko.
\newblock MLIR: Scaling compiler infrastructure for domain-specific computation.
\newblock In \emph{CGO}, pages 2--14, 2021. DOI: 10.1109/CGO51591.2021.9370308.

\bibitem{llvmnpm}
LLVM Project.
\newblock Using the New Pass Manager.
\newblock Official LLVM documentation, accessed 2026.

\bibitem{mlirpm}
LLVM Project.
\newblock MLIR Pass Management.
\newblock Official MLIR documentation, accessed 2026.

\bibitem{pnueli98}
Amir Pnueli, Michael Siegel, and Eli Singerman.
\newblock Translation validation.
\newblock In \emph{TACAS}, pages 151--166, 1998. DOI: 10.1007/BFb0054170.

\bibitem{necula00}
George C. Necula.
\newblock Translation validation for an optimizing compiler.
\newblock In \emph{PLDI}, pages 83--94, 2000. DOI: 10.1145/349299.349314.

\bibitem{tristan08}
Jean-Baptiste Tristan and Xavier Leroy.
\newblock Formal verification of translation validators: A case study on instruction scheduling optimizations.
\newblock In \emph{POPL}, pages 17--27, 2008. DOI: 10.1145/1328438.1328444.

\bibitem{leroy06}
Xavier Leroy.
\newblock Formal certification of a compiler back-end, or: Programming a compiler with a proof assistant.
\newblock In \emph{POPL}, pages 42--54, 2006. DOI: 10.1145/1111037.1111042.

\bibitem{leroy09backend}
Xavier Leroy.
\newblock A formally verified compiler back-end.
\newblock \emph{Journal of Automated Reasoning}, 43(4):363--446, 2009. DOI: 10.1007/s10817-009-9155-4.

\bibitem{alive}
Nuno P. Lopes, David Menendez, Santosh Nagarakatte, and John Regehr.
\newblock Provably correct peephole optimizations with Alive.
\newblock In \emph{PLDI}, pages 22--32, 2015. DOI: 10.1145/2737924.2737965.

\bibitem{alive2}
Nuno P. Lopes, Juneyoung Lee, Chung-Kil Hur, Zhengyang Liu, and John Regehr.
\newblock Alive2: Bounded translation validation for LLVM.
\newblock In \emph{PLDI}, pages 65--79, 2021. DOI: 10.1145/3453483.3454030.

\bibitem{pcc}
George C. Necula.
\newblock Proof-carrying code.
\newblock In \emph{POPL}, pages 106--119, 1997. DOI: 10.1145/263699.263712.

\bibitem{tal}
Greg Morrisett, David Walker, Karl Crary, and Neal Glew.
\newblock From System F to typed assembly language.
\newblock \emph{ACM TOPLAS}, 21(3):527--568, 1999. DOI: 10.1145/319301.319345.

\bibitem{certcompiler}
George C. Necula and Peter Lee.
\newblock The design and implementation of a certifying compiler.
\newblock In \emph{PLDI}, pages 333--344, 1998. DOI: 10.1145/277652.277752.

\bibitem{csmith}
Xuejun Yang, Yang Chen, Eric Eide, and John Regehr.
\newblock Finding and understanding bugs in C compilers.
\newblock In \emph{PLDI}, pages 283--294, 2011. DOI: 10.1145/1993498.1993532.

\bibitem{emi}
Vu Le, Mehrdad Afshari, and Zhendong Su.
\newblock Compiler validation via equivalence modulo inputs.
\newblock In \emph{PLDI}, pages 216--226, 2014. DOI: 10.1145/2594291.2594334.

\bibitem{yarpgen}
Vsevolod Livinskii, Dmitry Babokin, and John Regehr.
\newblock Random testing for C and C++ compilers with YARPGen.
\newblock \emph{PACMPL}, 4(OOPSLA), 2020. DOI: 10.1145/3428264.

\bibitem{creduce}
John Regehr, Yang Chen, Pascal Cuoq, Eric Eide, Chucky Ellison, and Xuejun Yang.
\newblock Test-case reduction for C compiler bugs.
\newblock In \emph{PLDI}, pages 335--346, 2012. DOI: 10.1145/2254064.2254104.

\bibitem{perses}
Chengnian Sun, Yuanbo Li, Qirun Zhang, Tianxiao Gu, and Zhendong Su.
\newblock Perses: Syntax-guided program reduction.
\newblock In \emph{ICSE}, pages 361--371, 2018. DOI: 10.1145/3180155.3180236.

\bibitem{delta}
Andreas Zeller and Ralf Hildebrandt.
\newblock Simplifying and isolating failure-inducing input.
\newblock \emph{IEEE Transactions on Software Engineering}, 28(2):183--200, 2002. DOI: 10.1109/32.988498.

\end{thebibliography}

\end{document}
